\section{Exploratory Robustness Results}

The extension attempted 8,960 replicate triads, generating 26,880 trajectories. The common observation gate failed for 968 triads, all in combined adversity; each generator retained 39 evaluable triads from 160 attempts in that arm. Failure requires any one of the three independent role trajectories to lack the minimum observations. These results therefore cannot be pooled with the previous single-trajectory failure rate. The complete archive is registered as \texttt{M02-DEV-ROBUSTNESS-2026-10-08}.

There were 663 candidate-level rank failures, distinct from failed triads. Seven-category observation made intercept/AR actions constant, leaving 1,920 full-rank candidate fits without a finite Gaussian variance MLE and 1,280 mode/criterion records without a selection. These records count as unsuccessful recovery; their unavailable AIC/BIC values are not replaced by an arbitrary variance floor. The latter count repeats choices across modes/criteria and is not an additional number of failed triads.

\begin{table}[htbp]
\centering
\caption{Exploratory nominal recovery under selection rules}
\label{tab:m02-robustness-selection}
\begin{tabular}{llrrrr}
\toprule
Arm & Generator & Temporal & Independent & AIC & BIC \\
\midrule
Reference & P & 26/160 & 31/160 & 97/160 & 147/160 \\
Reference & PID & 149/160 & 150/160 & 160/160 & 160/160 \\
Noise & P & 9/160 & 7/160 & 30/160 & 83/160 \\
Noise & PID & 83/160 & 85/160 & 82/160 & 34/160 \\
\bottomrule
\end{tabular}
\end{table}


Table~\ref{tab:m02-robustness-selection} compares nominal recovery. Temporal and independent columns use validation RMSE; AIC/BIC use the same fitting data and consequently have identical recovery across evaluation modes by construction. In the reference arm, independent validation selected P in 31/160 attempts versus 147/160 for BIC, while PID recovery was 150/160 versus 160/160. Under Gaussian analyst noise, BIC selected P in 83/160 attempts but PID in only 34/160, compared with 7/160 and 85/160 for independent validation RMSE. Penalization can resolve some nested-family labeling ambiguity while suppressing a noisy higher-order generator. These paired developmental counts do not establish a universally superior selector or a statistically established gain from independent validation.

Independent validation-RMSE PID recovery was 64/160 under variance-matched Student-$t$ noise, 32/160 with irregular sampling, 60/160 with seven-category quantization and 40/160 under clipping, versus 150/160 in the reference arm. Combined adversity yielded 5/160 attempts, or 5/39 evaluable triads. Non-Gaussian noise and altered observation schedules thus remain relevant stressors in this specified design. Quantized-response RMSE is on category-center units; its predictive performance cannot be read as unaltered latent-action performance.

For reference PID, the independent validation-RMSE mean paired loss difference was approximately $-0.000861$, with percentile bootstrap interval $[-0.000936,-0.000779]$. It remains within the original development tolerance. With combined adversity, the corresponding interval was approximately $[-0.010660,0.007407]$, crossing both equivalence boundaries. A bootstrap does not turn sparse, unstable extrapolation into a precise result.

\begin{table}[htbp]
\centering
\caption{Factorial sensitivity: pointwise bootstrap labels}
\label{tab:m02-robustness-epsilon}
\begin{tabular}{rrrrr}
\toprule
$\epsilon$ & Equivalent & Inconclusive & Generator & Competitor \\
\midrule
0.001 & 168 & 199 & 99 & 46 \\
0.0025 & 313 & 134 & 48 & 17 \\
0.005 & 404 & 69 & 24 & 15 \\
0.01 & 472 & 29 & 4 & 7 \\
\bottomrule
\end{tabular}
\end{table}


The factorial reanalysis classified 404/512 cells as development-equivalent at $\epsilon=0.005$ under percentile bootstrap, compared with 405/512 under the archived normal approximation. Eleven cell labels changed at that fixed tolerance across all categories. At $\epsilon=0.001$, only 168/512 were inside the band; at 0.01, 472/512 were inside. Table~\ref{tab:m02-robustness-epsilon} reports all four thresholds, including inconclusive and both directions of beyond-band separation. The overall similarity between interval methods at 0.005 does not establish coverage, and strong tolerance sensitivity precludes a threshold-independent equivalence claim.

PID fitting matrices were full rank in all evaluable extension triads. Median conditioning after slope-feature normalization was approximately 3.81 in the reference arm, 3.75 with Gaussian analyst noise and 7.48 in combined adversity. Full rank and modest conditional conditioning can coexist with failed family recovery: local precision conditional on noisy reconstructed inputs is not accurate recovery of latent gains. The indexed working profiles diagnose a fixed-design regression only. They do not distinguish structural non-identifiability of the closed-loop generator from measurement-induced bias.

The extension and factorial reanalysis retain their exploratory status. Fresh seeds and untouched final scoring strengthen the evaluation design, but do not undo prior knowledge of the development findings or provide a new preregistered confirmation.
